\chapter{Markov chain Monte Carlo}
\label{ch:mcmc}

\section{Sampling a posterior}

The goal of Bayesian calibration is the posterior $p(\theta \mid y)$ of \cref{eq:bayes} --- not
one best point, but the full distribution of parameters consistent with the data: their means,
widths, correlations, and credible intervals. In more than a few dimensions a grid over parameter
space is hopeless, because its size grows exponentially with $d$. Markov chain Monte Carlo (MCMC)
instead builds a random walk whose steps are accepted or rejected so that, in the long run, the
points it visits are distributed according to $p$. Histograms of those points then estimate the
posterior, and averages over them estimate posterior expectations. Crucially, a step only ever
compares $p$ at two points, so the unknown normalizing constant cancels.

\section{The affine-invariant ensemble sampler}

slurmcmc uses the ensemble sampler of \texttt{emcee}. Instead of one chain it advances an ensemble
of $\Nc$ \textbf{walkers} --- $\Nc$ parallel MCMC chains, whose starting points are
\code{init_points}. A walker $\theta_k$ moves by the \emph{stretch move}: pick another walker
$\theta_j$ and propose a point on the line through both,
\begin{equation}
  \theta' = \theta_j + z\,(\theta_k - \theta_j),
  \qquad z \sim g(z) \propto \frac{1}{\sqrt{z}}\ \ \text{on}\ \ \left[\tfrac{1}{a},\, a\right],
\end{equation}
with $a = 2$, and accept it with probability
\begin{equation}
  \min\!\left(1,\; z^{\,d-1}\,\frac{p(\theta')}{p(\theta_k)}\right).
\end{equation}
Because proposals are built from the positions of the other walkers, the sampler adapts
automatically to the scale and orientation of the posterior: it is invariant under any linear
transformation of the parameters, so strongly correlated, elongated posteriors are sampled as
easily as round ones.

\begin{keypoint}[Why the ensemble parallelizes]
The walkers are split into two halves, and each half is updated using the positions of the other.
All proposals of one half can therefore be evaluated \emph{at the same time}. slurmcmc submits each
half as one job array, so an iteration costs two arrays of $\Nc/2$ evaluations, preceded once by an
array evaluating all $\Nc$ starting points. The parallel speed-up is therefore bounded by $\Nc/2$
simultaneous evaluations, and \texttt{emcee} requires at least $\Nc \ge 2d$ walkers.
\end{keypoint}

\section{When has a chain converged?}
\label{sec:mcmc-convergence}

Successive points of a chain are correlated, so $\Nc \Lc$ samples --- $\Nc$ walkers times $\Lc$
iterations --- carry much less information than as many independent draws. Convergence diagnostics
quantify how much, following the \href{https://emcee.readthedocs.io/en/stable/tutorials/autocorr/}{\texttt{emcee}
documentation}.

\subsection{Integrated autocorrelation time}

For a parameter with normalized autocorrelation function $\rho(t)$ along the chain, the integrated
autocorrelation time is
\begin{equation}
  \tau = \sum_{t=-\infty}^{\infty} \rho(t) = 1 + 2 \sum_{t=1}^{\infty} \rho(t),
\end{equation}
roughly the number of iterations between effectively independent samples. It is computed per
parameter, using all walkers together to reduce the noise of the estimate. The estimate itself is
only trustworthy when the chain is much longer than $\tau$:
\begin{equation}
  \Lc \ge 50\,\tau .
  \label{eq:lc-tau}
\end{equation}

\begin{cautionbox}[$\tau$ grows before it settles]
The estimated $\tau$ usually \emph{increases} with $\Lc$ until it saturates, because a short chain
cannot see its own slowest correlations. $\Lc/\tau$ therefore does not simply grow in proportion
to $\Lc$, and \cref{eq:lc-tau} is only meaningful once $\tau$ has stopped growing. More walkers
improve the estimate of $\tau$, but do not make $\tau$ converge at a shorter $\Lc$.
\end{cautionbox}

\subsection{Effective sample size}

The effective sample size is the number of independent draws the chains are worth,
\begin{equation}
  \ESS = \frac{N_{\mathrm{total}}}{\tau} = \frac{\Nc \Lc}{\tau},
  \label{eq:ess}
\end{equation}
advisably 10 to 100 times the number of parameters. Unlike $\Lc/\tau$, it grows in proportion to
the number of walkers --- a wide ensemble can show a healthy ESS while each chain is far too short,
which is why both criteria are needed.

\subsection{Supplementary checks}

\begin{itemize}
  \item \textbf{Gelman--Rubin statistic.} The ratio of between-chain to within-chain variance,
        $\hat R$, approaches 1 at convergence. Walkers of an ensemble sampler are correlated by
        construction, so it is only a supplementary diagnostic here.
  \item \textbf{Stopping points.} Plot the posterior as it would have been had the chain stopped at
        several points along the way. Agreement between the last few is necessary, though not
        sufficient, for convergence.
  \item \textbf{$\chi^2$ distribution.} When the posterior is close to Gaussian,
        $-2\log p$ (with the normalization of the example) follows a $\chi^2$ distribution with
        $d$ degrees of freedom, whose mean is $d$. Comparing the histogram of accepted values with
        it is a quick sanity check.
\end{itemize}

The first part of every chain, while the walkers travel from their starting points into the
posterior, is not representative and is discarded as \textbf{burn-in}; a few autocorrelation
times, or a fixed fraction of the chain, is typical.

\needspace{16\baselineskip}
\section{Usage}

\begin{lstlisting}
import numpy as np
from scipy.optimize import rosen
from slurmcmc.mcmc import slurm_mcmc

def log_prob(theta):
    return -rosen(theta)                       # up to a constant

num_params, num_walkers = 2, 10
x_min = np.array([1.0, 1.0])                   # e.g. from slurm_minimize
init_points = x_min + 0.1 * np.random.randn(num_walkers, num_params)

status = slurm_mcmc(log_prob_fun=log_prob, init_points=init_points, num_iters=10000,
                    cluster='slurm', work_dir='mcmc', save_restart=True)

sampler = status['sampler']                    # an emcee.EnsembleSampler
tau = sampler.get_autocorr_time(quiet=True)
samples = sampler.get_chain(discard=int(2 * tau.max()), flat=True)
\end{lstlisting}

The status dictionary holds the \code{emcee} \code{sampler} --- whose \code{get_chain},
\code{get_log_prob} and \code{get_autocorr_time} give everything needed for the analysis --- the
\code{slurm_pool} with the evaluation history, the number of the next iteration \code{ini_iter},
and \code{time_per_iter}. The sampler advances one iteration at a time, so a restart file can be
written after every iteration.

\begin{argtable}
  \code{log_prob_fun} & required & the log-probability, up to a constant; a callable or a deferred-import dictionary \\
  \code{init_points} & required & starting walkers, shape $(\Nc, d)$, with $\Nc \ge 2d$ \\
  \code{num_iters} & required & iterations to run (further iterations, when resuming) \\
  \code{init_log_prob_fun_values} & \code{None} & log-probabilities already known at \code{init_points}, to skip their evaluation \\
  \code{emcee_kwargs} & \code{None} & forwarded to \code{emcee.EnsembleSampler}, e.g. \code{moves} \\
  \code{record_history} & \code{True} & keep the evaluation history; \code{False} saves memory on very long runs \\
  \code{progress}, \code{print_iter_interval} & \code{False}, \code{1} & \texttt{emcee}'s progress bar, and how often iterations are logged \\
  \code{skip_initial_state_check} & \code{True} & passed to \texttt{emcee} \\
  \code{status_restart} & \code{None} & resume from a status dictionary in memory \\
  \code{job_fail_value} & $-10^{10}$ & log-probability recorded for a failed evaluation \\
\end{argtable}

\begin{keypoint}[Constraints and starting points]
\code{slurm_mcmc} has no separate constraint argument: return a very low log-probability outside
the allowed region, and the walkers will not enter it. Start the walkers inside that region, and
preferably close to the optimum found in \cref{ch:optimization}, so the burn-in is short.
\end{keypoint}

\section{Example: a constrained 2D Rosenbrock}
\label{sec:mcmc-example}

The log-probability is minus the two-dimensional Rosenbrock function, with the circular constraint
of \cref{sec:opt-example} imposed by returning $-10^{7}$ outside the circle. $\Nc = 10$ walkers run
for $\Lc = 10^4$ iterations from random starting points inside the circle. For an expensive model in
higher dimension, the starting points should instead come from an optimization.

\begin{figure}[H]
  \centering
  \includegraphics[width=0.78\linewidth]{example_mcmc_chains_progress.png}
  \caption{The walkers' positions per iteration, after burn-in.}
\end{figure}

\begin{figure}[H]
  \centering
  \includegraphics[width=0.78\linewidth]{example_mcmc_convergence_diagnostics.png}
  \caption{Convergence diagnostics per parameter as functions of the chain length $\Lc$:
           the autocorrelation time and $\Lc/\tau$, the effective sample size, and the
           Gelman--Rubin statistic.}
\end{figure}

\begin{figure}[H]
  \centering
  \includegraphics[width=0.72\linewidth]{example_mcmc_2d_visualization.png}
  \caption{The log-probability (logarithm of its absolute value), with every point the walkers
           proposed (black) and the points accepted into the chains (red).}
\end{figure}

\begin{figure}[H]
  \centering
  \includegraphics[width=\linewidth]{example_mcmc_log_probability_values_explore.png}
  \caption{The evolution of $\chi^2 = -2\log p$ during the run, and the histogram of accepted
           values (red) against a $\chi^2$ distribution with two degrees of freedom (blue). The
           numerically computed means are given in the legend; the analytic mean of the $\chi^2$
           distribution is the dimension.}
\end{figure}

The final product of the analysis is a \href{https://github.com/dfm/corner.py}{\texttt{corner}}
plot of the posterior, and its stability along the chain:

\begin{figure}[H]
  \centering
  \includegraphics[width=0.62\linewidth]{example_mcmc_parameters_distribution.png}
  \caption{The posterior: one- and two-dimensional marginals.}
\end{figure}

\begin{figure}[H]
  \centering
  \includegraphics[width=0.62\linewidth]{example_mcmc_parameters_distribution_convergence.png}
  \caption{The posterior as it would have been had the chains been stopped at five points along
           the run, from purple (earliest) to red (the full chain).}
\end{figure}

\section{MCMC against grid Monte Carlo}
\label{sec:mcmc-vs-mc}

How much does MCMC actually save? The target is minus the five-dimensional Rosenbrock function of
\cref{eq:rosenbrock}, sampled with $\Nc = 50$ walkers for $\Lc = 5\cdot 10^4$ iterations, starting
close to the known minimum.

\begin{figure}[H]
  \centering
  \includegraphics[width=0.68\linewidth]{example_mcmc_5d_parameters_distribution_convergence.png}
  \caption{The 5D posterior at five stopping points along the chains, without burn-in.}
\end{figure}

\begin{figure}[H]
  \centering
  \includegraphics[width=0.62\linewidth]{example_mcmc_5d_convergence_diagnostics.png}
  \caption{Convergence diagnostics of the 5D run, as in \cref{sec:mcmc-convergence}.}
\end{figure}

The same posterior is then computed by brute force, evaluating the density on a regular grid:

\begin{figure}[H]
  \centering
  \includegraphics[width=0.7\linewidth]{example_mcmc_and_mc_5d_parameters_distribution.png}
  \caption{The 5D posterior from MCMC (black) and from grid Monte Carlo (blue), and from a
           coarser grid with about as many samples as the MCMC (red). The legend gives the total
           number of samples and the computation time of each.}
\end{figure}

\begin{keypoint}[What the comparison shows]
The fine grid needs about $40\times$ more evaluations than MCMC for the same distribution, although
its total time was only about $2.5\times$ longer here, because MCMC carries more algorithmic
overhead per sample. Given the \emph{same} number of samples, the coarse grid is far from converged.
The gap widens rapidly with dimension --- and when each evaluation is an expensive simulation, it is
the number of evaluations, not the overhead, that counts.
\end{keypoint}
